\begin{abstract}
Patch-based Transformer forecasters such as PatchTST treat every input patch as equally trustworthy, yet sensor series routinely contain missing readings, measurement noise, spikes and shifts in level or scale. This paper presents RA-PatchTST, a reliability-aware extension of PatchTST that adds three small components to the original design: a patch reliability gate (a two-layer perceptron with 801 parameters that scores each patch between 0 and 1 and rescales its embedding), a shift-aware forecasting head that receives the mean and scale of each input window relative to the training statistics, which instance normalization would otherwise discard, and corruption-aware training, in which half of the training windows are corrupted on the fly while the loss is always computed against the clean future. The full model adds 993 parameters (1.2\%) to an 81,728-parameter backbone. We evaluate on the ETTh1 benchmark (17,420 hourly records, seven variables, chronological 12/4/4-month split, look-back 336, horizon 96) using five ablation variants, two random seeds and 22 fixed test conditions: clean data plus 21 seeded corruptions covering point-wise and 16-step missing values, Gaussian noise, spikes, segment-level shift and a mixed setting. Our PatchTST re-implementation reaches a clean test MSE of 0.379 against the published 0.375. RA-PatchTST lowers the mean MSE over the 21 corrupted conditions from 0.429 to 0.413 ($-3.7\%$), with the largest reductions under shift ($-11.0\%$), spikes ($-6.5\%$) and mixed damage ($-5.2\%$), at a cost of $+0.7\%$ on clean data and up to $+1.3\%$ under missing values. The ablation attributes about four fifths of the gain to corruption-aware training; the gate and the shift-aware head add a further $-0.8\%$, concentrated on shift. The gate assigns clearly lower weight to corrupted patches only for spikes. With one dataset, one horizon, two seeds and synthetic corruptions drawn from the training families, these results are preliminary evidence for the approach and do not establish general superiority.
\end{abstract}

\begin{IEEEkeywords}
Time-series forecasting, PatchTST, Transformer, robustness, distribution shift, data corruption, reliability gating, instance normalization, ablation study.
\end{IEEEkeywords}
